\sectionpaged{Pressure}
\index{Atmospheric pressure}
\index{Pressure!Atmospheric}

Atmospheric pressure\sidesupercite{WikiAtmosphericPressure}\index{Atmospheric pressure}\index{Pressure!Atmospheric},
also known as air pressure or barometric pressure (after the barometer),
is the pressure within the atmosphere of Earth.

\marginnote{
    \begin{kaobox}[title=Equivalent Units]%
        \begin{eqnarray*}%
            \Atmospheres{1} &\equiv& \hPascal{1,013.25} \\
            &\equiv& \Pascal{101,325} \\
            &\equiv& \Millibars{1,013.25} \\
            &\equiv& \mmMercury{760} \\
            &\equiv& \inchesMercury{29.9212} \\
            &\equiv& \PSI{14.696}
        \end{eqnarray*}
    \end{kaobox}
}The standard atmosphere (symbol: \Atmospheres{}) is a unit of pressure defined as
\Pascal{101,325} (\hPascal{1,013.25}).
The \Atmospheres{} unit is roughly equivalent to the mean sea-level atmospheric pressure on Earth;
that is, the Earth's atmospheric pressure at sea level is approximately \Atmospheres{1}.

\subsection[Mean sea-level]{Mean sea-level pressure}
\label{subsec:mean-sea-level-pressure}
\index{Mean Sea level}\index{Pressure!Mean Sea level}
Mean sea level\sidesupercite{WikiMeanSeaLevel} (MSL, often shortened to sea level)
is an average surface level of one or more among Earth's coastal bodies of water from
which heights such as elevation may be measured.
The global MSL is a type of vertical datum – a standardised geodetic datum – that is used,
for example, as a chart datum in cartography and marine navigation,
or, in aviation, as the standard sea level at which atmospheric pressure
is measured to calibrate altitude and, consequently, aircraft flight levels.
A common and relatively straightforward mean sea-level standard is instead a
long-term average of tide gauge readings at a particular reference location.
